\chapter*{Lời Mở Đầu}
\addcontentsline{toc}{chapter}{Lời Mở Đầu}
\markboth{Lời Mở Đầu}{Lời Mở Đầu}

\begin{truyen}{Lời Mở Đầu}{Opening Words}
\chuhoa{T}{uổi trẻ hay nuôi dưỡng ảo tưởng — đó không phải tội lỗi, nhưng cũng không phải điều đáng giữ mãi.}
\textit{(Youth tends to nurture illusions --- that is not a sin, but it is not something worth keeping forever.)}

Ảo tưởng về năng lực của mình, ảo tưởng về sự công bằng của thế giới, ảo tưởng rằng mọi thứ sẽ tự nhiên tốt hơn mà không cần làm gì. Quyển mười này gọi tên những ảo tưởng phổ biến nhất của người trẻ.
\textit{(Illusions about one's own abilities, illusions about the fairness of the world, illusions that everything will naturally get better without effort. Volume ten names the most common illusions of young people.)}

Không phải để phán xét, mà để nhận ra và điều chỉnh sớm hơn.
\textit{(Not to judge, but to recognize and correct them sooner.)}

Nhận ra ảo tưởng là bước đầu tiên để bước vào thực tế một cách chủ động hơn.
\textit{(Recognizing an illusion is the first step toward engaging with reality more proactively.)}
\end{truyen}

\begin{baihoc}
Ảo tưởng không xấu --- nhưng sống mãi trong ảo tưởng là cách tốt nhất để bỏ lỡ cuộc đời thật.
\textit{(Illusions are not bad --- but living in them forever is the best way to miss your real life.)}
\end{baihoc}

\begin{ghinhoanh}
\textbf{Short English to remember:}

Illusions are comfortable but costly.

Seeing clearly is an act of courage.
\end{ghinhoanh}

\begin{tuhocnhanh}
1. Em từng nuôi một ảo tưởng nào và sau đó nhận ra nó không thật? \textit{Have you ever held an illusion that you later realized was not real?}

2. Điều gì giúp em nhận ra ảo tưởng của mình? \textit{What helped you recognize your illusion?}

3. Sau khi đọc quyển này, em muốn điều chỉnh điều gì trong cách nhìn về bản thân và thế giới? \textit{After reading this volume, what do you want to adjust in how you see yourself and the world?}
\end{tuhocnhanh}

\ngancach
